\documentclass{article}

% if you need to pass options to natbib, use, e.g.:
%     \PassOptionsToPackage{numbers, compress}{natbib}
% before loading neurips_2026

% The authors should use one of these tracks.
% Before accepting by the NeurIPS conference, select one of the options below.
% 0. "default" for submission


% Compilation vars
\newif\ifreview \newcommand{\review}{\reviewtrue}
\newif\ifarxiv \newcommand{\arxiv}{\arxivtrue}
\newif\ifcamera \newcommand{\cameraready}{\cameratrue}
\newif\ifrebuttal \newcommand{\rebuttal}{\rebuttaltrue}

\PassOptionsToPackage{numbers, compress}{natbib}

\review % \review OR \arxiv OR \cameraready
% \arxiv
% \rebuttal
%\cameraready

\ifreview \usepackage[main]{neurips_2026} \fi
\ifarxiv \usepackage[preprint]{neurips_2026} \fi
\ifrebuttal \usepackage[main]{neurips_2026} \fi
\ifcamera \usepackage[main, final]{neurips_2026} \fi

% the "default" option is equal to the "main" option, which is used for the Main Track with double-blind reviewing.
% 1. "main" option is used for the Main Track
%  \usepackage[main]{neurips_2026}
% 2. "position" option is used for the Position Paper Track
%  \usepackage[position]{neurips_2026}
% 3. "eandd" option is used for the Evaluations & Datasets Track
 % \usepackage[eandd]{neurips_2026}
 % if you need to opt-in for a single-blind submission in the E&D track:
 %\usepackage[eandd, nonanonymous]{neurips_2026}
% 4. "creativeai" option is used for the Creative AI Track
%  \usepackage[creativeai]{neurips_2026}
% 5. "sglblindworkshop" option is used for the Workshop with single-blind reviewing
 % \usepackage[sglblindworkshop]{neurips_2026}
% 6. "dblblindworkshop" option is used for the Workshop with double-blind reviewing
%  \usepackage[dblblindworkshop]{neurips_2026}

% After being accepted, the authors should add "final" behind the track to compile a camera-ready version.
% 1. Main Track
 % \usepackage[main, final]{neurips_2026}
% 2. Position Paper Track
%  \usepackage[position, final]{neurips_2026}
% 3. Evaluations & Datasets Track
 % \usepackage[eandd, final]{neurips_2026}
% 4. Creative AI Track
%  \usepackage[creativeai, final]{neurips_2026}
% 5. Workshop with single-blind reviewing
%  \usepackage[sglblindworkshop, final]{neurips_2026}
% 6. Workshop with double-blind reviewing
%  \usepackage[dblblindworkshop, final]{neurips_2026}
% Note. For the workshop paper template, both \title{} and \workshoptitle{} are required, with the former indicating the paper title shown in the title and the latter indicating the workshop title displayed in the footnote.
% For workshops (5., 6.), the authors should add the name of the workshop, "\workshoptitle" command is used to set the workshop title.
% \workshoptitle{WORKSHOP TITLE}

% "preprint" option is used for arXiv or other preprint submissions
 % \usepackage[preprint]{neurips_2026}

% to avoid loading the natbib package, add option nonatbib:
%    \usepackage[nonatbib]{neurips_2026}

\usepackage[utf8]{inputenc} % allow utf-8 input
\usepackage[T1]{fontenc}    % use 8-bit T1 fonts
\usepackage{hyperref}       % hyperlinks
\usepackage{url}            % simple URL typesetting
\usepackage{booktabs}       % professional-quality tables
\usepackage{amsfonts}       % blackboard math symbols
\usepackage{nicefrac}       % compact symbols for 1/2, etc.
\usepackage{microtype}      % microtypography
\usepackage{xcolor}         % colors

% Note. For the workshop paper template, both \title{} and \workshoptitle{} are required, with the former indicating the paper title shown in the title and the latter indicating the workshop title displayed in the footnote. 




\def\paperTitle{
% Turbo-Muon: Accelerating Orthogonality with Pre-Conditioning
% Turbo-Muon: Accelerating Orthogonality-Based Optimization with Pre-Conditioning
Turbo-Muon: Almost-Orthogonal Pre-Conditioning for Fast Muon Updates
%Turbo-Muon: When Newton-Schulz Preconditioning Leads to Improved Muon Runtime
% Turbo-Muon: Turbocharging Muon with AOL Pre-Conditioning
% Turbo-Muon/Polar-steps/flash-newton schulz: Accelerating orthogonality-based optimizer with pre-conditioning
}

\def\authorBlock{
\hspace{0cm}
\textbf{Thibaut Boissin}$^{1,2,\dagger}$
\hspace{1mm}
\textbf{Thomas Massena}$^{2,3,\dagger}$
\hspace{1mm}
\textbf{Franck Mamalet}$^{1}$
\hspace{1mm}
\textbf{Mathieu Serrurier}$^{2}$
\vspace{3mm}
\\
$^1$ Institut de Recherche Technologique Saint-Exupery, France\\
$^2$ IRIT, France\\
$^3$ Innovation \& Research Division, SNCF, France\\
    % Author 1\thanks{Equal contribution} \qquad
    % Author 2\footnotemark[1] \qquad
    % Author 3 \\
    % Institute \\
    % {\tt\small \{email, addresses\}@inst.edu}
}

%\title{Turbo-Muon: Accelerating Orthogonality-Based Optimization with Pre-Conditioning}
% Turbo-Muon: Accelerating Orthogonality with Pre-Conditioning

% Turbo-Muon: Turbocharging Muon with AOL Pre-Conditioning
% Turbo-Muon/Polar-steps/flash-newton schulz: Accelerating orthogonality-based optimizer with pre-conditioning


% The \author macro works with any number of authors. There are two commands
% used to separate the names and addresses of multiple authors: \And and \AND.
%
% Using \And between authors leaves it to LaTeX to determine where to break the
% lines. Using \AND forces a line break at that point. So, if LaTeX puts 3 of 4
% authors names on the first line, and the last on the second line, try using
% \AND instead of \And before the third author name.


%\author{%
%  Thibaut Boissin$^{1,3,\dagger}$ \\
%  $^1$ Institut de Recherche Technologique Saint-Exupery, France\\
%  \And
%  Thomas Massena$^{2,3,\dagger}$ \\
%    $^2$ Innovation \& Research Division, SNCF, France\\
%  \And
%  Franck Mamalet$^{1}$
%  \And
%  Mathieu Serrurier$^{3}$\\
%  $^3$ IRIT, France
  % examples of more authors
  % \And
  % Coauthor \\
  % Affiliation \\
  % Address \\
  % \texttt{email} \\
  % \AND
  % Coauthor \\
  % Affiliation \\
  % Address \\
  % \texttt{email} \\
  % \And
  % Coauthor \\
  % Affiliation \\
  % Address \\
  % \texttt{email} \\
  % \And
  % Coauthor \\
  % Affiliation \\
  % Address \\
  % \texttt{email} \\
%}


\usepackage{graphicx}	
\usepackage{amsmath}	
\usepackage{amssymb}	
\usepackage{amsthm}
\usepackage{booktabs}
\usepackage{times}
\usepackage{microtype}
\usepackage{epsfig}
\usepackage{caption}
\usepackage{float}
\usepackage{color, colortbl}
\usepackage{stfloats}
\usepackage{enumitem}

\usepackage{tabularx}
\usepackage{xstring}
\usepackage{multirow}
\usepackage{xspace}
\usepackage{url}
\usepackage{subcaption}
\usepackage{xcolor}
\usepackage[hang,flushmargin]{footmisc}
\usepackage{bbm}
%\usepackage[section]{placeins}
\usepackage{ulem}

\newtheorem{assumption}{Assumption}
\newtheorem{proposition}{Proposition}
\newtheorem{definition}{Definition}
\newtheorem{lemma}{Lemma}
\newtheorem{corollary}{Corollary}
\newtheorem{theorem}{Theorem}
\newtheorem{remark}{Remark}
%% For cross-referencing labels between documents
\usepackage{xr-hyper}
\usepackage{makecell}


\usepackage{wrapfig}

% Unfortunately, this package interferes with arxiv's stamp
\ifcamera \usepackage[accsupp]{axessibility} \fi

%% MACROS

% \newcommand{\authorname}[1]{{\textcolor{blue}{[Author: #1]}}}
% ...

% \newcommand{\commandname}{string\xspace}
% \definecolor{colorname}{rgb}{0.92,0.49,0.19}

\definecolor{muoncolor}{HTML}{0072B2}
\definecolor{dioncolor}{HTML}{56B4E9}
\definecolor{turbomuoncolor}{HTML}{009E73}
\definecolor{frobcolor}{HTML}{7E7E7E}
% General

\newcommand{\nbf}[1]{{\noindent \textbf{#1.}}}

\newcommand{\supp}{supplemental material\xspace}
\ifarxiv \renewcommand{\supp}{appendix\xspace} \fi

\newcommand{\todo}[1]{{\textcolor{red}{[TODO: #1]}}}

% Reviewer commands (1 to 5), e.g. \R{1}, \R{2}
\newcommand{\R}[1]{{%
    \textbf{%
        \ifstrequal{#1}{1}{\textcolor{red}{R#1}}{%
        \ifstrequal{#1}{2}{\textcolor{blue}{R#1}}{%
        \ifstrequal{#1}{3}{\textcolor{magenta}{R#1}}{%
        \ifstrequal{#1}{4}{\textcolor{teal}{R#1}}{%
                           \textcolor{cyan}{R#1}%
        }}}}%
    }%
}}


\newcommand{\turbomuon}{Turbo-Muon }

\newif\ifshowcomments
\showcommentstrue % Comment this line to disable comments

% Define the comment commands
\ifshowcomments

\newcommand{\thomas}[1]{{\color{blue}Thomas:#1}}
\newcommand{\thib}[1]{{\color{magenta}Thib:#1}}
\newcommand{\franck}[1]{{\color{red}Franck:#1}}
\else
\newcommand{\thib}[1]{}
\newcommand{\thomas}[1]{}
\newcommand{\franck}[1]{{}
\fi

% \widowpenalty=50     % default is 150, lower it to allow widows
% \clubpenalty=50      % default is 150, lower it to allow orphans
% \brokenpenalty=50    % default is 1000, allows breaks at hyphenated lines
% \displaywidowpenalty=50  % default is 50, allows widows in displayed equations
% \predisplaypenalty=50    % default is 10000, allows page breaks before displayed equations
% \setlength{\parskip}{-1pt}      % remove extra space between paragraphs
% % \setlength{\parskip}{1pt}      % remove extra space between paragraphs
% \setlength{\parindent}{1em}    % reduce paragraph indentation if necessary


\makeatletter
\newcommand*{\addFileDependency}[1]{
  \typeout{(#1)}
  \@addtofilelist{#1}
  \IfFileExists{#1}{}{\typeout{No file #1.}}
}

\makeatother
\newcommand*{\myexternaldocument}[1]{
    \externaldocument{#1}
    \addFileDependency{#1.tex}
    \addFileDependency{#1.aux}
}
%%

%\definecolor{cvprblue}{rgb}{0.21,0.49,0.74}
%\usepackage[pagebackref,breaklinks,colorlinks,allcolors=cvprblue]{hyperref}
\usepackage[capitalize]{cleveref}
\crefname{section}{Sec.}{Secs.}
\crefname{table}{Table}{Tables}
\crefname{figure}{Fig.}{Figs.}
\crefname{algorithm}{Alg.}{Algs.}

\ifarxiv \crefname{appendix}{App.}{Apps.}
\else \crefname{appendix}{Suppl.}{Suppls.} \fi

%\frenchspacing

\usepackage{algorithm}%
\usepackage{algorithmicx}%
\usepackage{algpseudocode}%

\newtheorem*{lemma*}{Lemma}